\subsection{浸没与商流形}

5.9.1. 设 \(f: X \to Y\) 是流形的一个态射，\(a\) 是 \(X\) 的一点，并置 \(b = f(a)\)。下列条件等价：

\begin{enumerate}
\def\labelenumi{(\roman{enumi})}
\item
  线性映射 \(T_{a}(f)\) 是满射，且其核在 \(T_{a}(X)\) 中具有拓扑补空间。
\item
  存在 X 在 \(\varphi\) 处的一个图 (U, \(a\) , E)、Y 在 \(\psi\) 处的一个图 (V, \(b\) , F)，以及从 E 到 F 的满连续线性映射 \(u\)，使得
\end{enumerate}

\[
f (\mathrm{U}) \subset \mathrm{V}, \quad \psi \circ f = u \circ \varphi
\]

且 \(u\) 的核在 E 中具有拓扑补空间。

\begin{enumerate}
\def\labelenumi{(\roman{enumi})}
\setcounter{enumi}{2}
\item
  存在 a 的一个开邻域 U、包含 \(f(\mathbf{U})\) 的 b 的一个开邻域 V，以及从 U 到流形 Z 的态射 g，使得映射 \((f,g)\) 从 U 到 \(V \times Z\) 是一个同构。
\item
  存在 b 的一个开邻域 V 和从 V 到 X 的态射 s，使得 \(s(b) = a\)，且对 V 中所有 y，都有 \(f(s(y)) = y\)（“局部截面”）。
\end{enumerate}

当 X 和 Y 都是有限维时，前述条件等价于下列条件：

\begin{enumerate}
\def\labelenumi{(\alph{enumi})}
\setcounter{enumi}{21}
\tightlist
\item
  存在 a 的一个开邻域 U、包含 \(f(\mathbf{U})\) 的 b 的一个开邻域 V，以及 V 上的一个坐标系 \((\eta^{1},\ldots,\eta^{n})\)，使得函数 \(\eta^{i}\circ f\) 在 U 上构成 U 上某个坐标系的一部分。
\end{enumerate}

若性质 (i) 至 (iv) 成立，则称 \(f\) 在 \(a\) 处是浸没的。\(f\) 为浸没的点集是 \(X\) 中的开集；若这个开集就是 X 的全部，则称 f 是浸没。

5.9.2. 设 \(f: X \to Y\) 和 \(g: Y \to Z\) 是两个态射。若 \(f\) 和 \(g\) 都是浸没，则 \(g \circ f\) 也是浸没。反之，若 \(g \circ f\) 是浸没且 \(f\) 是满射，则 \(g\) 是浸没。

5.9.3. 若 \(f: X \to Y\) 和 \(f': X' \to Y'\) 是浸没，则 \(f \times f'\) 是浸没。

5.9.4. 浸没 \(f: X \to Y\) 是开映射（参见一般拓扑，第 I 章，第 4 \(^{e}\) 版，§ 5，第 1 号）；特别地，由 \(f\) 定义的等价关系 R 是开的，f 通过取商定义了从 X/R 到 \(f\)\(f(X)\) 的同胚，且 \(f(X)\) 在 Y 中是开的。

5.9.5. 设 R 是流形 X 上的一个等价关系。若在商空间 X/R 上存在流形结构，使典范投影 \(p: X \to X/R\) 为浸没，则称 R 是正规的；此流形结构于是是唯一的；它是 X 的流形结构的商（集合论，第~IV~章，§ 2，第 6 号）：换言之，从 X/R 到流形 Y 的态射 \(g\) 存在的充分必要条件是 \(g \circ p\) 是从 X 到 Y 的态射。

令 \(C \subset X \times X\) 为 R 的图。R 为正规的充分必要条件是满足下列两个条件：

\begin{enumerate}
\def\labelenumi{(\roman{enumi})}
\item
  C 是 X × X 的子流形。
\item
  映射 \(pr_{1}\) 从 C 到 X 是浸没。
\end{enumerate}

(ii) 还意味着：若 \(a\) 和 \(b\) 模 \(\mathbf{R}\) 同余，则存在 a 的一个开邻域 U 和从 U 到 X 的态射 \(a\)\(s\)，使得 \(s(a) = b\)，且对所有 ，\(s(x)\) 模  与 \(x\) 同余。\(\mathbf{R}\)\(x \in \mathbf{U}\)

假设 R 是正规的。商流形 X/R 是分离的充分必要条件是 R 的图在 \(X \times X\) 中闭。

5.9.6. 设 X 和 X' 是两个流形，R 和 R' 分别是 X 和 X' 上的正规等价关系，\(f: X \to X'\) 是一个与关系 R 和 R' 相容的态射。由 f 通过取商得到的映射 \(\bar{f}: X/R \to X'/R'\) 是一个态射。\(f\)

5.9.7.（“商的传递性”）设 R 和 S 是流形 X 上的两个等价关系，且 R 蕴含 S；令 S/R 为 X/R 上的商等价关系。假设 R 是正规的。则 S 为正规的充分必要条件是 S/R 为正规；若条件成立，则典范双射

\[
(\mathrm{X} / \mathrm{R}) / (\mathrm{S} / \mathrm{R}) \rightarrow \mathrm{X} / \mathrm{S}
\]

是流形的同构。

5.9.8.（“商的乘积”）设 \((\mathbf{X}_{i})_{i\in I}\) 是有限个流形组成的族，每个流形都赋有一个正规等价关系 \(R_{i}\)。令 \(X = \prod_{i\in I} X_{i}\)，并令 R 为 X 上由 \(R_{i}\) 的乘积得到的等价关系（参见一般拓扑，第 I 章，第 4 版，§ 5，第 3 号，命题 8 的推论）。则 R 是正规的，且从 X/R 到 \(\prod_{i\in I}(X_{i}/R_{i})\) 的典范双射是流形同构。

